\subsection{Public Deliberation and Participatory AI Governance}
\label{sec:8.2.2}
What is specific to governance, as against development, is the machinery for bringing public input to bear on a decision at a single point in time instead of through an ongoing relationship.

\begin{enumerate}
  \item Citizen assemblies. A panel of citizens selected by random, stratified sampling deliberates on a question over an extended period and produces recommendations grounded in that deliberation rather than in a snap poll. Ireland's Citizens' Assembly has used exactly this format to shape national recommendations on abortion law and on climate policy — questions carrying the same mix of technical complexity and moral stakes that AI governance has \autocite{citizensassembly2016ireland}. No assembly has yet been convened on AI itself, which is a fact about political attention and not about the method.
  \item Online platforms for public input. A standing online consultation reaches people a physical assembly cannot, at the cost of the depth that extended facilitated deliberation provides. The European Commission's consultation on its AI White Paper drew more than 1,200 responses from citizens, companies, and civil-society groups, and fed into the legislative proposal that became the EU AI Act \autocite{europeancommission2020whitepaper}.
  \item AI ethics committees. An ethics board answerable only to the company that created it is only as independent as that company is willing to let it be. DeepMind Health set up an outside panel in 2016 to review its work with the UK's National Health Service and publish annual reports on its findings; Google disbanded the panel in 2019, after members had raised concerns about how much access to information, and how much independence from the company, they actually had \autocite{fisher2019deepmind}. The panel's short life is the governance lesson, and it is the answer to the question section~\ref{sec:8.2.1} leaves open: the board said something the company did not want to hear, and the company's available response was to stop having a board. Nothing in the arrangement had made that response costly. A review body is worth what it can enforce, and this one could enforce publication of its own concerns right up until the moment it could not.
\end{enumerate}

\runin{What none of these mechanisms does}
There is a limit these mechanisms share, and it took me most of this book to see it, because it is invisible as long as you are looking at them one at a time.

Every governance mechanism proposed in this chapter — citizen juries, standing consultations, participatory design, deliberative polling, the vTaiwan model — sits on the same side of a line. All of them make public preference more legible, more informed, and more actionable. Not one of them protects anybody from it. They are aggregation instruments, and a better aggregation instrument produces a more faithful rendering of what the public wants, including when what the public wants is something done to a minority of it.

This is not an objection to any of them. Preference aggregation is a real requirement and these are good ways to meet it. It is an observation that meeting it completely would leave the central problem untouched. A perfectly deliberative, perfectly informed, perfectly representative process that concludes a group should be surveilled has produced a legitimate output by every criterion in this chapter.

Section~\ref{sec:7.2} gets close to this. It notices that treating all disagreement as equally worth representing goes wrong when one side's position is the product of coercion, which is correct. But it frames the trouble as a data-quality problem — the input is contaminated, so clean the input. The structural version of the claim is stronger and does not depend on anyone's preference being contaminated at all: some protections have to hold against the aggregate regardless of how the aggregate came out, and no amount of improvement to the aggregating machinery can supply them.

What supplies them is a different kind of machinery, and this chapter contains none of it: enforceable rights that do not depend on being popular, a court willing to rule against a government that holds a legitimate majority, limits that whoever won the last election cannot repeal, and standing for people who are affected without being enfranchised. Section~\ref{sec:10.10} takes up why that half matters most exactly where a democracy is behaving most democratically, and why its absence is what let this book's account of democratic accountability look complete when it was half a theory.
